\subsection{Mémoire généalogique et contenu informationnel d'une lecture}
\label{sec:ii-memoria-lectura}

L’opérateur d’aire réunit des occupations de frontière. L’observable
pondérée ajoute leur provenance sectorielle. Cette différence exprime un
principe physique du manuscrit : l’information accessible dépend de
l’application qui extrait la coordonnée observable de l’état. Dans APP, le résidu et le quotient
conservent conjointement l’opération entière ; dans TPK, la phase visible
coexiste avec l’orientation et l’accroissement de mémoire. La réalisation
physique transporte cette même distinction entre état, observable et
registre de récupération \cite{informacion}.

Soit \(\Omega\) un ensemble fini d’états atteints par la
mise à jour considérée. Un lecteur \(q:\Omega\to Y\) assigne la
coordonnée observable et \(M:\Omega\to\mathcal M\) conserve la mémoire
de la fibre : feuillet, orientation, quotient, retenue et chemin pertinents.
Pour une distribution \(p\) sur ces états, notons
\(\mathsf H_p\) l’entropie adimensionnelle, avec le logarithme naturel.

\begin{proposition}[Résolution informationnelle de la fibre]
Si \((q,M)\) est injective sur le support de \(p\), alors
\begin{align}
 \mathsf H_p(X)&=\mathsf H_p(q(X))+
                 \mathsf H_p(X\mid q(X)),\\
 \mathsf H_p(X\mid q(X),M(X))&=0.
 \label{eq:ii-memoria-condicional}
\end{align}
\end{proposition}
\begin{proof}
La factorisation de la distribution conjointe en distribution visible
et distribution conditionnelle sur la fibre donne la première égalité.
L’injectivité fournit un inverse sur l’image de \((q,M)\) ;
chaque paire exprimée détermine un unique état du support et sa
distribution conditionnelle est concentrée sur cet état.
\end{proof}

L’équation \eqref{eq:recover} réalise ce mécanisme pour les deux
occupations agrégées : \(N\) donne le total et \(D\) résout leur
distribution sectorielle. Les états \((N_{90},N_{120})=(1,0)\) et
\((0,1)\) partagent l’aire agrégée et se distinguent par la seconde
observable. Sous une distribution équiprobable sur cette paire, la
lecture exclusive de \(N\) laisse \(\log2\) nats d’incertitude ;
la paire \((N,D)\) la résout exactement. L’affirmation concerne
ces occupations. Les microhistoires au sein de chaque secteur conservent
leur registre supplémentaire.

L'interprétation de la mémoire de bord est ainsi déterminée par
une récupération concrète. Une réduction peut masquer une information
que l'état conjoint continue de transporter. Pour une mise à jour
bijective \(U\), l'égalité \(\mathsf H_p(X\mid U(X))=0\) exprime
la réversibilité logique. Une opération physique d'effacement ou de réinitialisation
de la mémoire possède, quant à elle, son propre bilan d'information et
d'énergie. L'intrication s'évalue relativement à la bipartition et à
l'état réduit spécifiés dans la section précédente.

\subsection{Feuillets aréal et volumétrique : incidence, croissance et mesure}
\label{sec:ii-hojas-av}

La distinction aréal--volumétrique possède une réalisation dynamique dans le
calendrier du TPK. Le bloc \((+1,-1,0)\), appliqué aux modes
\(\Sigma,\Pi\) par
\[
 +1:m\mapsto\Sigma,\qquad -1:m\mapsto\Pi,\qquad 0:m\mapsto m,
\]
produit l’orbite cyclique \((\Sigma,\Pi,\Pi)\). Le lecteur géométrique
réalise les correspondances \(\Sigma\mapsto\mathscr H_{\rm ar}\) et
\(\Pi\mapsto\mathscr H_{\rm vol}\). Leurs transitions consécutives
fixent l’incidence
\begin{equation}
 F_{\rm av}=\begin{pmatrix}0&1\\1&1\end{pmatrix},\qquad
 N_9=3F_{\rm av},\quad N_{27}=9F_{\rm av},\quad
 N_{108}=36F_{\rm av}.
 \label{eq:ii-hojas-incidencia}
\end{equation}
En effet, le cycle contient une transition de chaque type
\(\Sigma\to\Pi\), \(\Pi\to\Pi\), \(\Pi\to\Sigma\).
Les calendriers plus grands répètent ce bloc avec les multiplicités
indiquées.

L'enveloppe symbolique associée permet toutes les concaténations
admises par \(F_{\rm av}\). Son nombre de chemins dépend des
puissances de cette matrice. Avec \(F_0=0\), \(F_1=1\), la récurrence donne
\begin{equation}
 F_{\rm av}^{\,n}=
 \begin{pmatrix}F_{n-1}&F_n\\F_n&F_{n+1}\end{pmatrix},
 \qquad h_{\rm top}=\log\varphi.
 \label{eq:ii-hojas-fibonacci}
\end{equation}
L'identité matricielle utilise la récurrence de Fibonacci. La croissance
exponentielle est fixée par la racine de Perron \(\varphi\) de
l'incidence déjà construite.

\begin{proposition}[Distribution stationnaire des feuillets]
La mesure d'entropie maximale de l'enveloppe symbolique a pour poids
\begin{equation}
 \mu_{\rm ar}=\frac{1}{1+\varphi^2},\qquad
 \mu_{\rm vol}=\frac{\varphi^2}{1+\varphi^2},\qquad
 \frac{\mu_{\rm ar}}{\mu_{\rm vol}}=\varphi^{-2}.
 \label{eq:ii-hojas-medida}
\end{equation}
\end{proposition}
\begin{proof}
Le vecteur positif \(r=(1,\varphi)\) satisfait
\(F_{\rm av}r=\varphi r\). Les probabilités de transition
\(P_{ij}=(F_{\rm av})_{ij}r_j/(\varphi r_i)\) ont pour somme un.
La symétrie de \(F_{\rm av}\) donne des poids stationnaires proportionnels
à \(r_i^2\). Leur normalisation produit
\eqref{eq:ii-hojas-medida}. Dans la moyenne stationnaire,
\(-\log P_{ij}=\log\varphi+\log r_i-\log r_j\) sur les
arêtes permises. Les deux derniers termes s'annulent et
l'entropie par transition est \(\log\varphi\), égale à l'entropie
topologique de \eqref{eq:ii-hojas-fibonacci}.
\end{proof}

Le calendrier périodique et son enveloppe symbolique constituent deux
objets explicites : le premier possède les fréquences chronologiques
\((1/3,2/3)\) ; le second admet les concaténations de l’incidence
et possède la mesure \eqref{eq:ii-hojas-medida}. Cette séparation permet
d’attribuer à chaque quantité son sens physico-mathématique. Le rapport
\(\varphi^{-2}\) caractérise la distribution des feuillets dans cette
mesure ; \(\log\varphi\) caractérise la croissance de ses histoires
permises. La lecture angulaire de frontière incorpore ensuite la
coordonnée régionale correspondante. La construction distingue
croissance volumétrique, représentation aréale et mémoire du transport.

\subsection{Continuations admissibles et information de survie}
\label{sec:ii-informacion-supervivencia}

Un état TPK \(x\) étant fixé, soit \(\Omega_n(x)\) l’ensemble fini
de ses continuations de longueur \(n\) admises par la règle de
prolongement. Les restrictions entre niveaux conservent les préfixes.
Une famille compatible de probabilités \(p_n\), obtenue à partir du
lecteur statistique déclaré, définit l’incertitude de
continuation
\begin{equation}
 s_n(x)=-\sum_{\omega\in\Omega_n(x)}
                 p_n(\omega)\log p_n(\omega).
 \label{eq:ii-entropia-continuaciones}
\end{equation}
La compatibilité exige que la probabilité d’un préfixe soit la somme
des probabilités de ses prolongements. Dans l’enveloppe précédente,
la mesure stationnaire fournit une famille de ce type.

\begin{proposition}[Bilan informationnel du raffinement]
Pour la famille compatible précédente,
\begin{align}
 0\leq s_n(x)&\leq\log|\operatorname{supp}p_n|,\\
 s_{n+1}(x)-s_n(x)
 &=\mathsf H(X_{n+1}\mid X_n)\geq0.
 \label{eq:ii-entropia-refinamiento}
\end{align}
La borne supérieure est atteinte pour la distribution uniforme sur le
support. L'incrément est nul lorsque chaque préfixe de probabilité
positive détermine son prolongement avec probabilité un.
\end{proposition}
\begin{proof}
L'inégalité élémentaire pour l'entropie d'une distribution finie
donne la première ligne. Puisque \(X_n\) est la restriction de \(X_{n+1}\),
la factorisation des probabilités sur chaque fibre de prolongement
donne \(\mathsf H(X_{n+1})=\mathsf H(X_n)+
\mathsf H(X_{n+1}\mid X_n)\). Chaque entropie conditionnelle est non
négative et s'annule exactement pour une distribution ponctuelle.
\end{proof}

Les continuations qui demeurent admissibles possèdent ainsi un contenu
informationnel précis. Leur multiplicité et leurs poids décrivent
l'incertitude résiduelle de la lecture à profondeur finie. Une
histoire réalisée conserve, en outre, son orientation et sa mémoire.
L'entropie informationnelle, l'entropie de l'état réduit et l'entropie
thermodynamique se relient par leurs lecteurs respectifs ; chacune
conserve l'information de domaine nécessaire à son interprétation.
